\ifdefined\gkver\else\def\gkver{student}\fi
\documentclass[\gkver]{gaokaozhenti}
\begin{document}

\chapter{2012年全国理综}
%% paperType: 理综物理部分

\begin{enumerate}
\item
%% number: 14
%% paperName: 2012年全国理综第 14 题 6 分
%% typeId: 2
%% type: 多选题
%% chapter: 气体性质
%% point: 分子动理论
%% method: 无
%% score: 6
%% degree: 650
%% duplicateId: 0
%% body:
下列关于布朗运动的说法，正确的是\xzanswer{BD}
\fourchoices[answer=BD]
{布朗运动是液体分子的无规则运动}
{液体温度越高，悬浮粒子越小，布朗运动越剧烈}
{布朗运动是由于液体各部分的温度不同而引起的}
{布朗运动是由液体分子从各个方向对悬浮粒子撞击作用的不平衡引起的}
%% answer: BD
%% memo:
\memoanswer{
布朗运动是悬浮颗粒的无规则运动，不是液体分子的运动，选项 A 错误；液体温度越高，悬浮颗粒越小，布朗运动越剧烈，选项 B 正确；布朗运动是由于液体分子从各个方向对悬浮粒子撞击作用不平衡引起的，选项 C 错误，选项 D 正确。

正确选项：BD。
}

\item
%% number: 15
%% paperName: 2012年全国理综第 15 题 6 分
%% typeId: 1
%% type: 单选题
%% chapter: 原子和原子核
%% point: 天然放射性
%% method: 无
%% score: 6
%% degree: 650
%% duplicateId: 0
%% body:
$^{235}_{92}$U经过$m$次$\alpha$衰变和$n$次$\beta$衰变变为$^{207}_{82}$Pb，则\xzanswer{B}
\fourchoices[answer=B]
{$m=7$，$n=3$}
{$m=7$，$n=4$}
{$m=14$，$n=9$}
{$m=14$，$n=18$}
%% answer: B
%% memo:
\memoanswer{
原子核每发生一次 $\alpha$ 衰变，质量数减少 4，电荷数减少 2；每发生一次 $\beta$ 衰变，质量数不变，电荷数增加 1。比较两种原子核，质量数减少 28，即发生了 7 次 $\alpha$ 衰变；电荷数应减少 14，而实际电荷数减少 10，说明发生了 4 次 $\beta$ 衰变，故 B 项正确。

正确选项：B。
}

\item
%% number: 16
%% paperName: 2012年全国理综第 16 题 6 分
%% typeId: 2
%% type: 多选题
%% chapter: 光学
%% point: 光的衍射
%% method: 无
%% score: 6
%% degree: 650
%% duplicateId: 0
%% body:
在双缝干涉实验中，某同学用黄光作为入射光，为了增大干涉条纹的间距，该同学可以采用的方法有\xzanswer{AC}
\fourchoices[answer=AC]
{改用红光作为入射光}
{改用蓝光作为入射光}
{增大双缝到屏的距离}
{增大双缝之间的距离}
%% answer: AC
%% memo:
\memoanswer{
根据 $\Delta x=\frac{L}{d}\lambda$，红光比黄光波长大，所以 A 正确；蓝光比黄光波长短，所以 B 错误；增大双缝到屏的距离即增大 $L$，增大干涉条纹的间距，所以 C 正确；增大双缝之间的距离，减小干涉条纹的间距，所以 D 错误。

正确选项：AC。
}

\item
%% number: 17
%% paperName: 2012年全国理综第 17 题 6 分
%% typeId: 1
%% type: 单选题
%% chapter: 磁场
%% point: 带电粒子在磁场中的运动
%% method: 无
%% score: 6
%% degree: 650
%% duplicateId: 0
%% body:
质量分别为$m_{1}$和$m_{2}$、电荷量分别为$q_{1}$和$q_{2}$的两粒子在同一匀强磁场中做匀速圆周运动，已知两粒子的动量大小相等。下列说法正确的是\xzanswer{A}
\fourchoices[answer=A]
{若$q_{1}=q_{2}$，则它们作圆周运动的半径一定相等}
{若$m_{1}=m_{2}$，则它们作圆周运动的半径一定相等}
{若$q_{1}\neq q_{2}$，则它们作圆周运动的周期一定不相等}
{若$m_{1}\neq m_{2}$，则它们作圆周运动的周期一定不相等}
%% answer: A
%% memo:
\memoanswer{
根据牛顿第二定律得 $qvB=m\frac{v^{2}}{r}$，由圆周运动规律得 $T=\frac{2\pi r}{v}$。由此得到轨道半径 $r=\frac{mv}{qB}=\frac{p}{qB}$ 及周期 $T=\frac{2\pi m}{qB}$。

由于两粒子动量 $p$ 相等，若 $q_{1}=q_{2}$，则半径 $r$ 一定相等，A 正确；半径还与电荷量有关，B 错误；周期只与 $m$、$q$ 有关，$q$ 不同时周期不一定不同，$m$ 不同时周期也不一定不同，C、D 错误。

正确选项：A。
}

\item
%% number: 18
%% paperName: 2012年全国理综第 18 题 6 分
%% typeId: 1
%% type: 单选题
%% chapter: 磁场
%% point: 磁场的叠加
%% method: 无
%% score: 6
%% degree: 450
%% duplicateId: 0
%% body:
如图，两根互相平行的长直导线过纸面上的 M、N 两点，且与纸面垂直，导线中通有大小相等、方向相反的电流。a、O、b 在 M、N 的连线上，O 为 MN 的中点，c、d 位于 MN 的中垂线上，且 a、b、c、d 到 O 点的距离均相等。关于以上几点处的磁场，下列说法正确的是\xzanswer{C}
\onepicture{figs/18.png}
\fourchoices[answer=C]
{O 点处的磁感应强度为零}
{a、b 两点处的磁感应强度大小相等，方向相反}
{c、d 两点处的磁感应强度大小相等，方向相同}
{a、c 两点处磁感应强度的方向不同}
%% answer: C
%% memo:
\memoanswer{
由安培定则可知，两导线在 O 点产生的磁场均竖直向下，合磁感应强度一定不为零，选项 A 错；

由安培定则，两导线在 a、b 两处产生磁场方向均竖直向下，由于对称性，电流 M 在 a 处产生磁场的磁感应强度等于电流 N 在 b 处产生磁场的磁感应强度，同时电流 M 在 b 处产生磁场的磁感应强度等于电流 N 在 a 处产生磁场的磁感应强度，所以 a、b 两处磁感应强度大小相等方向相同，选项 B 错；

根据安培定则，两导线在 c、d 处产生磁场垂直 c、d 两点与导线连线方向向下，且产生的磁场的磁感应强度相等，由平行四边形定则可知，c、d 两点处的磁感应强度大小相同，方向相同，选项 C 正确。

a、c 两处磁感应强度的方向均竖直向下，选项 D 错。

正确选项：C。
}

\item
%% number: 19
%% paperName: 2012年全国理综第 19 题 6 分
%% typeId: 1
%% type: 单选题
%% chapter: 交流电
%% point: 交流电
%% method: 无
%% score: 6
%% degree: 450
%% duplicateId: 0
%% body:
一台电风扇的额定电压为交流$220\UV$。在其正常工作过程中，用交流电流表测得某一段时间内的工作电流$I$随时间$t$的变化如图所示。这段时间内电风扇的用电量为\xzanswer{B}
\onepicture{figs/19.png}
\fourchoices[answer=B]
{$3.9\times10^{-4}$度}
{$5.5\times10^{-2}$度}
{$7.8\times10^{-2}$度}
{$11.0\times10^{-2}$度}
%% answer: B
%% memo:
\memoanswer{
根据电流的变化情况，分段计算电功（时间单位换算为秒）：

$W_{1}=I_{1}Ut_{1}=0.3\UA\times220\UV\times600\Us=3.96\times10^{4}\UJ$；

$W_{2}=I_{2}Ut_{2}=0.4\UA\times220\UV\times600\Us=5.28\times10^{4}\UJ$；

$W_{3}=I_{3}Ut_{3}=0.2\UA\times220\UV\times2400\Us=1.056\times10^{5}\UJ$。

则总功 $W=W_{1}+W_{2}+W_{3}=1.98\times10^{5}\UJ=5.5\times10^{-2}$ 度。

正确选项：B。
}

\item
%% number: 20
%% paperName: 2012年全国理综第 20 题 6 分
%% typeId: 2
%% type: 多选题
%% chapter: 机械振动
%% point: 振动图象
%% method: 无
%% score: 6
%% degree: 450
%% duplicateId: 0
%% body:
一列简谐横波沿$x$轴正方向传播，图（a）是$t=0$时刻的波形图，图（b）和图（c）分别是$x$轴上某两处质点的振动图像。由此可知，这两质点平衡位置之间的距离可能是\xzanswer{BD}
\onepicture{figs/20.png}
\fourchoices[answer=BD]
{$\frac{1}{3}\Um$}
{$\frac{2}{3}\Um$}
{$1\Um$}
{$\frac{4}{3}\Um$}
%% answer: BD
%% memo:
\memoanswer{
若 b 在前 c 在后，则 b、c 间相差为 $\frac{1}{3}T$，这两质点平衡位置之间的距离可能是 $\frac{1}{3}\lambda=\frac{2}{3}\Um$（$\lambda=2\Um$），B 正确。

若 b 在后 c 在前，则 c、b 间相差为 $\frac{2}{3}T$，这两质点平衡位置之间的距离可能是 $\frac{2}{3}\lambda=\frac{4}{3}\Um$（$\lambda=2\Um$），D 正确。

正确选项：BD。
}

\item
%% number: 21
%% paperName: 2012年全国理综第 21 题 6 分
%% typeId: 2
%% type: 多选题
%% chapter: 运动和力
%% point: 动量守恒定律
%% method: 无
%% score: 6
%% degree: 450
%% duplicateId: 0
%% body:
如图，大小相同的摆球a和b的质量分别为$m$和$3m$，摆长相同，并排悬挂，平衡时两球刚好接触，现将摆球a向左边拉开一小角度后释放，若两球的碰撞是弹性的，下列判断正确的是\xzanswer{AD}
\onepicture{figs/21.png}
\fourchoices[answer=AD]
{第一次碰撞后的瞬间，两球的速度大小相等}
{第一次碰撞后的瞬间，两球的动量大小相等}
{第一次碰撞后，两球的最大摆角不相同}
{发生第二次碰撞时，两球在各自的平衡位置}
%% answer: AD
%% memo:
\memoanswer{
根据碰撞过程动量守恒和动能守恒，得 $m_{1}v_{1}=m_{1}v_{1}'+m_{2}v_{2}'$，$\frac{1}{2}m_{1}v_{1}^{2}=\frac{1}{2}m_{1}{v_{1}'}^{2}+\frac{1}{2}m_{2}{v_{2}'}^{2}$，且 $m_{1}=m$，$m_{2}=3m$，解得 $v_{1}'=\frac{m_{1}-m_{2}}{m_{1}+m_{2}}v_{1}=-\frac{1}{2}v_{1}$，$v_{2}'=\frac{2m_{1}}{m_{1}+m_{2}}v_{1}=\frac{1}{2}v_{1}$，所以 A 正确，B 错误；

根据 $\frac{1}{2}mv^{2}=mgh=mgR(1-\cos\theta)$，知第一次碰撞后两球的最大摆角 $\theta$ 相同，C 错误；根据单摆的等时性，D 正确。

正确选项：AD。
}

\item
%% number: 22
%% paperName: 2012年全国理综第 22 题 6 分
%% typeId: 4
%% type: 实验题
%% chapter: 电路
%% point: 电路实验
%% method: 无
%% score: 6
%% degree: 650
%% duplicateId: 0
%% body:
在黑箱内有一由四个阻值相同的电阻构成的串并联电路，黑箱面板上有三个接线柱1、2和3。用欧姆表测得1、2接线柱之间的电阻为$1\Omega$，2、3接线柱之间的电阻为$1.5\Omega$，1、3接线柱之间的电阻为$2.5\Omega$。
\onepicture[align=right]{figs/22.png}
\begin{enumerate}
	\item
	在虚线框中画出黑箱中的电阻连接方式；
	\item
	如果将1、3接线柱用导线连接起来，1、2接线柱之间的电阻为\_\_\_\_\_\_$\Omega$。
\end{enumerate}
%% answer: 
\jdanswer{
\begin{enumerate}
	\item
	如图。
	\item
	$0.6\Omega$
\end{enumerate}
}
%% memo:
\memoanswer{
\onepicture[h=3cm, align=right]{figs/22_an.png}

（1）因为 1、2 接线柱之间的电阻与 2、3 接线柱之间的电阻之和等于 1、3 接线柱之间的电阻，所以 2 为中间的结点；又因为 2、3 接线柱之间的电阻与 1、2 接线柱之间的电阻之差等于 1、2 接线柱之间的电阻的一半，故 2、3 之间有两个电阻并联，并联后再与第三个电阻串联，每个电阻均为 $1\Omega$，连接方式如图所示。

（2）将 1、3 用导线相连后，两个 $1\Omega$ 的电阻并联后与第三个 $1\Omega$ 的电阻串联，总电阻为 $1.5\Omega$，再与第四个 $1\Omega$ 的电阻并联，等效电阻 $R=0.6\Omega$。
}

\item
%% number: 23
%% paperName: 2012年全国理综第 23 题 11 分
%% typeId: 4
%% type: 实验题
%% chapter: 运动和力
%% point: 验证牛顿第二定律
%% method: 无
%% score: 11
%% degree: 450
%% duplicateId: 0
%% body:
图1为验证牛顿第二定律的实验装置示意图。图中打点计时器的电源为$50\UHz$的交流电源，打点的时间间隔用$\Delta t$表示。在小车质量未知的情况下，某同学设计了一种方法用来研究“在外力一定的条件下，物体的加速度与其质量间的关系”。
\begin{enumerate}
	\item
	完成下列实验步骤中的填空：

	①平衡小车所受的阻力：小吊盘中不放物块，调整木板右端的高度，用手轻拨小车，直到打点计时器打出一系列\_\_\_\_\_\_\_\_的点。

	②按住小车，在小吊盘中放入适当质量的物块，在小车中放入砝码。

	③打开打点计时器电源，释放小车，获得带有点列的纸带，在纸带上标出小车中砝码的质量$m$。

	④按住小车，改变小车中砝码的质量，重复步骤③。

	⑤在每条纸带上清晰的部分，每5个间隔标注一个计数点。测量相邻计数点的间距$s_{1}$，$s_{2}$，$\ldots$。求出与不同$m$相对应的加速度$a$。

	⑥以砝码的质量$m$为横坐标，$\frac{1}{a}$为纵坐标，在坐标纸上做出$\frac{1}{a}-m$关系图线。若加速度与小车和砝码的总质量成反比，则$\frac{1}{a}$与$m$应成\_\_\_\_\_\_\_\_\_关系（填“线性”或“非线性”）。

	\item
	完成下列填空：

	（i）本实验中，为了保证在改变小车中砝码的质量时，小车所受的拉力近似不变，小吊盘和盘中物块的质量之和应满足的条件是\_\_\_\_\_\_\_\_\_\_\_\_\_\_\_\_\_\_\_\_\_\_\_\_。

	（ii）设纸带上三个相邻计数点的间距为$s_{1}$、$s_{2}$和$s_{3}$。$a$可用$s_{1}$、$s_{3}$和$\Delta t$表示为$a=$\_\_\_\_\_\_\_\_\_\_。图2为用米尺测量某一纸带上的$s_{1}$、$s_{3}$的情况，由图可读出$s_{1}=$\_\_\_\_\_\_\_\_$\Umm$，$s_{3}=$\_\_\_\_\_\_\_\_$\Umm$。由此求得加速度的大小$a=$\_\_\_\_\_\_\_\_$\Umsq$。

	（iii）图3为所得实验图线的示意图。设图中直线的斜率为$k$，在纵轴上的截距为$b$，若牛顿定律成立，则小车受到的拉力为\_\_\_\_\_\_\_\_，小车的质量为\_\_\_\_\_\_\_\_。
\end{enumerate}
\twopicture[labelA=2012全国理综23a, labelB=2012全国理综23b, align=right]
{figs/23a.png}
{figs/23b.png}
%% answer: 
\jdanswer{
\begin{enumerate}
	\item
	间隔均匀；线性
	\item
	远小于小车的质量；$\frac{s_{3}-s_{1}}{50(\Delta t)^{2}}$；24.3；47.2；1.15；$\frac{1}{k}$；$\frac{b}{k}$
\end{enumerate}
}
%% memo:
\memoanswer{
（1）①平衡好小车所受的阻力，小车做匀速运动，打点计时器打出的点间隔基本相等。⑥根据牛顿第二定律可知 $F=(M+m)a$，得 $\frac{1}{a}=\frac{M}{F}+\frac{m}{F}$，$\frac{1}{a}$ 与 $m$ 为一次函数关系，是线性关系。

（2）（i）为保证小车所受拉力近似不变，应满足小吊盘和盘中物块的质量之和远小于小车的质量。

（ii）由 $\Delta x=aT^{2}$ 可知，$a=\frac{s_{3}-s_{1}}{2(5\Delta t)^{2}}=\frac{s_{3}-s_{1}}{50(\Delta t)^{2}}$。由图可读出 $s_{1}=36.8-12.5=24.3\Umm$，$s_{3}=120.0-72.8=47.2\Umm$，$a=\frac{s_{3}-s_{1}}{2(5\Delta t)^{2}}=\frac{(47.2-24.3)\times10^{-3}}{2\times(5\times0.02)^{2}}\Umsq=1.15\Umsq$。

（iii）设小车的质量为 $m'$，则有 $F=(m+m')a$，变形得 $\frac{1}{a}=\frac{1}{F}m+\frac{m'}{F}$，所以 $\frac{1}{a}-m$ 图象的斜率为 $\frac{1}{F}=k$，所以作用力 $F=\frac{1}{k}$，图象的截距为 $\frac{m'}{F}=b$，所以 $m'=\frac{b}{k}$。
}

\item
%% number: 24
%% paperName: 2012年全国理综第 24 题 16 分
%% typeId: 6
%% type: 计算题
%% chapter: 电路
%% point: 电容
%% method: 无
%% score: 16
%% degree: 450
%% duplicateId: 0
%% body:
如图，一平行板电容器的两个极板竖直放置，在两极板间有一带电小球，小球用一绝缘轻线悬挂于O点。现给电容器缓慢充电，使两极板所带电荷量分别为$+Q$和$-Q$，此时悬线与竖直方向的夹角为$\pi/6$。再给电容器缓慢充电，直到悬线和竖直方向的夹角增加到$\pi/3$，且小球与两极板不接触。求第二次充电使电容器正极板增加的电荷量。
\onepicture[align=right]{figs/24.png}
%% answer: 
\jdanswer{
$2Q$
}
%% memo:
\memoanswer{
设电容器的电容为 $C$，第一次充电后两极板间电压为 $U=\frac{Q}{C}$，两板之间电场的强度为 $E=\frac{U}{d}$，式中 $d$ 为两板间距离。

当小球偏转角 $\theta_{1}=\frac{\pi}{6}$ 时，小球处于平衡位置。设小球质量为 $m$，所带电荷量为 $q$，则有 $T\cos\theta_{1}=mg$，$T\sin\theta_{1}=qE$，式中 $T$ 为此时悬线的张力。联立以上各式得 $\tan\theta_{1}=\frac{qQ}{mgCd}$。

设第二次充电使正极板增加的电荷量为 $\Delta Q$，此时小球偏转角 $\theta_{2}=\frac{\pi}{3}$，则 $\tan\theta_{2}=\frac{q(Q+\Delta Q)}{mgCd}$。两式相除得 $\frac{\tan\theta_{1}}{\tan\theta_{2}}=\frac{Q}{Q+\Delta Q}$。

代入数据解得 $\Delta Q=2Q$。
}

\item
%% number: 25
%% paperName: 2012年全国理综第 25 题 19 分
%% typeId: 6
%% type: 计算题
%% chapter: 机械振动
%% point: 单摆
%% method: 无
%% score: 19
%% degree: 450
%% duplicateId: 0
%% body:
一单摆在地面处的摆动周期与在某矿井底部摆动周期的比值为$k$。设地球的半径为$R$。假定地球的密度均匀。已知质量均匀分布的球壳对壳内物体的引力为零，求矿井的深度$d$。
%% answer: 
\jdanswer{
$R(1-k^{2})$
}
%% memo:
\memoanswer{
根据万有引力定律，地面处质量为 $m$ 的物体的重力为 $mg=G\frac{mM}{R^{2}}$，式中 $g$ 是地面处的重力加速度，$M$ 是地球的质量。设 $\rho$ 是地球的密度，则有 $M=\frac{4}{3}\pi\rho R^{3}$。

摆长为 $L$ 的单摆在地面处的摆动周期为 $T=2\pi\sqrt{\frac{L}{g}}$。

若该物体位于矿井底部，则其重力为 $mg'=G\frac{mM'}{(R-d)^{2}}$，式中 $g'$ 是矿井底部的重力加速度，且 $M'=\frac{4}{3}\pi\rho(R-d)^{3}$。在矿井底部此单摆的周期为 $T'=2\pi\sqrt{\frac{L}{g'}}$。

由题意 $T=kT'$，联立以上各式得 $d=R(1-k^{2})$。
}

\item
%% number: 26
%% paperName: 2012年全国理综第 26 题 20 分
%% typeId: 6
%% type: 计算题
%% chapter: 曲线运动
%% point: 平抛运动
%% method: 无
%% score: 20
%% degree: 250
%% duplicateId: 0
%% body:
一探险队员在探险时遇到一山沟，山沟的一侧竖直，另一侧的坡面呈抛物线形状。此队员从山沟的竖直一侧，以速度$v_{0}$沿水平方向跳向另一侧坡面。如图所示，以沟底的$O$点为原点建立坐标系$Oxy$。已知，山沟竖直一侧的高度为$2h$，坡面的抛物线方程为$y=\frac{1}{2h}x^{2}$，探险队员的质量为$m$。人视为质点，忽略空气阻力，重力加速度为$g$。
\begin{enumerate}
	\item
	求此人落到坡面时的动能；
	\item
	此人水平跳出的速度为多大时，他落在坡面时的动能最小？动能的最小值为多少？
\end{enumerate}
\onepicture[align=right]{figs/26.png}
%% answer: 
\jdanswer{
\begin{enumerate}
	\item
	$E_{k}=\frac{1}{2}m\left(v_{0}^{2}+\frac{4g^{2}h^{2}}{v_{0}^{2}+gh}\right)$
	\item
	$\sqrt{gh}$，$\frac{3}{2}mgh$
\end{enumerate}
}
%% memo:
\memoanswer{
（1）设探险队员在空中运动的时间为 $t$，在坡面上落点的横坐标为 $x$、纵坐标为 $y$。由运动学公式和已知条件得 $x=v_{0}t$，$2h-y=\frac{1}{2}gt^{2}$，$y=\frac{1}{2h}x^{2}$。由机械能守恒，落到坡面时的动能为 $\frac{1}{2}mv^{2}=\frac{1}{2}mv_{0}^{2}+mg(2h-y)$。联立以上各式得 $\frac{1}{2}mv^{2}=\frac{1}{2}m\left(v_{0}^{2}+\frac{4g^{2}h^{2}}{v_{0}^{2}+gh}\right)$。

（2）上式可以改写为 $v^{2}=\left(\sqrt{v_{0}^{2}+gh}-\frac{2gh}{\sqrt{v_{0}^{2}+gh}}\right)^{2}+3gh$。$v^{2}$ 极小的条件为上式中的平方项等于 0，由此得 $v_{0}=\sqrt{gh}$。此时 $v^{2}=3gh$，则最小动能为 $\left(\frac{1}{2}mv^{2}\right)_{\min}=\frac{3}{2}mgh$。
}

\end{enumerate}

\end{document}
